%=============================================================================
%  Sunshine CTF 2025 pwn write-up — XeLaTeX template
%=============================================================================
\documentclass[10pt,a4paper]{article}

% ---------- 版面幾何 ----------
\usepackage[top=2.1cm,bottom=1.9cm,left=1.7cm,right=1.7cm,headsep=0.45cm,footskip=0.95cm]{geometry}

% ---------- 字型（XeLaTeX + xeCJK）----------
\usepackage{xeCJK}
\setCJKmainfont{Noto Sans CJK TC}
\setCJKsansfont{Noto Sans CJK TC}
\setCJKmonofont{Noto Sans Mono CJK TC}
\setmainfont{Noto Sans CJK TC}
\setsansfont{Noto Sans CJK TC}
\setmonofont[Scale=0.78]{Noto Sans Mono CJK TC}

% 中文斷行與標點
\xeCJKsetup{
  CheckSingle = false,
  CJKecglue = {\hskip 0.06em plus 0.02em},
}
\XeTeXlinebreaklocale "zh"
\XeTeXlinebreakskip = 0pt plus 1pt
\linespread{1.13}

% ---------- 顏色 ----------
\usepackage{xcolor}
\definecolor{dshblue}{HTML}{1F4E79}
\definecolor{dshcyan}{HTML}{1B6B7A}
\definecolor{dshgray}{HTML}{555555}
\definecolor{dshrule}{HTML}{B8C4D0}
\definecolor{codebg}{HTML}{F6F7F9}
\definecolor{codeframe}{HTML}{D8DEE6}
\definecolor{linkblue}{HTML}{0B5FA5}
\definecolor{quotebar}{HTML}{5B8DB8}

% ---------- 表格 ----------
\usepackage{longtable,booktabs,array}
\usepackage{etoolbox}
\renewcommand{\arraystretch}{1.28}
\setlength{\tabcolsep}{5pt}
\AtBeginEnvironment{longtable}{\small}
\AtBeginEnvironment{tabular}{\small}
\setlength{\LTpre}{0.7em}
\setlength{\LTpost}{0.7em}

% ---------- 程式碼區塊 ----------
\usepackage{fvextra}
\fvset{
  fontsize=\fontsize{7.1}{8.6}\selectfont,
  breaklines=true,
  breakanywhere=true,
  breakautoindent=false,
  frame=none,
  numbers=none,
  samepage=false,
}
\usepackage{framed}
\colorlet{shadecolor}{codebg}

% pandoc 產生的程式碼高亮巨集（必須先載入，才會定義 Shaded）
\usepackage{color}
\usepackage{fancyvrb}
\newcommand{\VerbBar}{|}
\newcommand{\VERB}{\Verb[commandchars=\\\{\}]}
\DefineVerbatimEnvironment{Highlighting}{Verbatim}{commandchars=\\\{\}}
% Add ',fontsize=\small' for more characters per line
\usepackage{framed}
\definecolor{shadecolor}{RGB}{248,248,248}
\newenvironment{Shaded}{\begin{snugshade}}{\end{snugshade}}
\newcommand{\AlertTok}[1]{\textcolor[rgb]{0.94,0.16,0.16}{#1}}
\newcommand{\AnnotationTok}[1]{\textcolor[rgb]{0.56,0.35,0.01}{\textbf{\textit{#1}}}}
\newcommand{\AttributeTok}[1]{\textcolor[rgb]{0.13,0.29,0.53}{#1}}
\newcommand{\BaseNTok}[1]{\textcolor[rgb]{0.00,0.00,0.81}{#1}}
\newcommand{\BuiltInTok}[1]{#1}
\newcommand{\CharTok}[1]{\textcolor[rgb]{0.31,0.60,0.02}{#1}}
\newcommand{\CommentTok}[1]{\textcolor[rgb]{0.56,0.35,0.01}{\textit{#1}}}
\newcommand{\CommentVarTok}[1]{\textcolor[rgb]{0.56,0.35,0.01}{\textbf{\textit{#1}}}}
\newcommand{\ConstantTok}[1]{\textcolor[rgb]{0.56,0.35,0.01}{#1}}
\newcommand{\ControlFlowTok}[1]{\textcolor[rgb]{0.13,0.29,0.53}{\textbf{#1}}}
\newcommand{\DataTypeTok}[1]{\textcolor[rgb]{0.13,0.29,0.53}{#1}}
\newcommand{\DecValTok}[1]{\textcolor[rgb]{0.00,0.00,0.81}{#1}}
\newcommand{\DocumentationTok}[1]{\textcolor[rgb]{0.56,0.35,0.01}{\textbf{\textit{#1}}}}
\newcommand{\ErrorTok}[1]{\textcolor[rgb]{0.64,0.00,0.00}{\textbf{#1}}}
\newcommand{\ExtensionTok}[1]{#1}
\newcommand{\FloatTok}[1]{\textcolor[rgb]{0.00,0.00,0.81}{#1}}
\newcommand{\FunctionTok}[1]{\textcolor[rgb]{0.13,0.29,0.53}{\textbf{#1}}}
\newcommand{\ImportTok}[1]{#1}
\newcommand{\InformationTok}[1]{\textcolor[rgb]{0.56,0.35,0.01}{\textbf{\textit{#1}}}}
\newcommand{\KeywordTok}[1]{\textcolor[rgb]{0.13,0.29,0.53}{\textbf{#1}}}
\newcommand{\NormalTok}[1]{#1}
\newcommand{\OperatorTok}[1]{\textcolor[rgb]{0.81,0.36,0.00}{\textbf{#1}}}
\newcommand{\OtherTok}[1]{\textcolor[rgb]{0.56,0.35,0.01}{#1}}
\newcommand{\PreprocessorTok}[1]{\textcolor[rgb]{0.56,0.35,0.01}{\textit{#1}}}
\newcommand{\RegionMarkerTok}[1]{#1}
\newcommand{\SpecialCharTok}[1]{\textcolor[rgb]{0.81,0.36,0.00}{\textbf{#1}}}
\newcommand{\SpecialStringTok}[1]{\textcolor[rgb]{0.31,0.60,0.02}{#1}}
\newcommand{\StringTok}[1]{\textcolor[rgb]{0.31,0.60,0.02}{#1}}
\newcommand{\VariableTok}[1]{\textcolor[rgb]{0.00,0.00,0.00}{#1}}
\newcommand{\VerbatimStringTok}[1]{\textcolor[rgb]{0.31,0.60,0.02}{#1}}
\newcommand{\WarningTok}[1]{\textcolor[rgb]{0.56,0.35,0.01}{\textbf{\textit{#1}}}}

\renewenvironment{Shaded}{%
  \def\FrameCommand{\fboxsep=3.2pt\colorbox{shadecolor}}%
  \MakeFramed{\advance\hsize-\width\FrameRestore\small}%
}{%
  \endMakeFramed
}

% ---------- 標題樣式 ----------
\usepackage{titlesec}
\titleformat{\section}{\normalfont\Large\bfseries\color{dshblue}}{\thesection}{0.6em}{}
\titleformat{\subsection}{\normalfont\large\bfseries\color{dshcyan}}{\thesubsection}{0.6em}{}
\titleformat{\subsubsection}{\normalfont\normalsize\bfseries\color{dshgray}}{\thesubsubsection}{0.6em}{}
\titleformat{\paragraph}[runin]{\normalfont\normalsize\bfseries}{}{0em}{}[~]
\titlespacing*{\section}{0pt}{1.6em}{0.55em}
\titlespacing*{\subsection}{0pt}{1.15em}{0.4em}
\titlespacing*{\subsubsection}{0pt}{0.9em}{0.3em}

% ---------- 頁眉頁腳 ----------
\usepackage{fancyhdr}
\pagestyle{fancy}
\fancyhf{}
\fancyhead[L]{\footnotesize\color{dshgray}Sunshine CTF 2025 — pwn 全 7 題 Write-up}
\fancyhead[R]{\footnotesize\color{dshgray}\leftmark}
\fancyfoot[C]{\footnotesize\color{dshgray}\thepage}
\renewcommand{\headrulewidth}{0.35pt}
\renewcommand{\footrulewidth}{0pt}
\renewcommand{\headrule}{\hbox to\headwidth{\color{dshrule}\leaders\hrule height \headrulewidth\hfill}}
\fancypagestyle{plain}{\fancyhf{}\fancyfoot[C]{\footnotesize\color{dshgray}\thepage}\renewcommand{\headrulewidth}{0pt}}
% 讓 \section 更新頁眉
\renewcommand{\sectionmark}[1]{\markboth{#1}{}}

% ---------- 引言區塊 ----------
\usepackage{mdframed}
\renewenvironment{quote}{%
  \begin{mdframed}[linewidth=0pt,innerleftmargin=0pt,innerrightmargin=0pt,
    innerbottommargin=0pt,innertopmargin=0pt,skipabove=0.55em,skipbelow=0.55em,
    leftmargin=0.9em,rightmargin=0.4em,startcode={\color{quotebar}\vrule width 2.6pt},%
    endcode={}]%
  \itshape\color{black}%
}{%
  \end{mdframed}%
}

% ---------- 清單 ----------
\usepackage{enumitem}
\setlist[itemize]{topsep=0.28em,itemsep=0.12em,parsep=0.1em,leftmargin=1.5em}
\setlist[enumerate]{topsep=0.28em,itemsep=0.12em,parsep=0.1em,leftmargin=1.7em}

% ---------- 連結 ----------
\usepackage[hidelinks,bookmarks=true,bookmarksnumbered=true]{hyperref}
\hypersetup{
  colorlinks=true,
  linkcolor=linkblue,
  urlcolor=linkblue,
  citecolor=linkblue,
  pdftitle={},
  pdfauthor={Sunshine CTF 2025 pwn write-up},
  pdfsubject={CTF pwn write-up},
  pdfcreator={pandoc + XeLaTeX},
  bookmarksopen=true,
  bookmarksopenlevel=1,
  breaklinks=true,
}

% ---------- 破折號、省略號等 ----------
\usepackage{amssymb}


% ---------- 標題 ----------
\title{\vspace{-1.2cm}\Huge\bfseries\color{dshblue} Sunshine CTF 2025\\[0.25em] pwn 全 7 題完整 Write-up\\[0.35em]
       \large\normalfont\color{dshgray} 從零開始：組語、漏洞成因、exploit 逐行解說}
\author{\normalsize 全部 7 題皆已對遠端 \texttt{chal.sunshinectf.games} 實打取得真實 flag}
\date{\normalsize\today}

\begin{document}
\sloppy
\emergencystretch=2.5em

{\hypersetup{linkcolor=dshblue}
\maketitle
\thispagestyle{plain}

\begin{center}
\begin{minipage}{0.9\textwidth}
\small\itshape\color{dshgray}
本文件是把 \texttt{wp/} 底下的教學文件合併成的單一 PDF。
每一題都包含：漏洞成因、逐行組合語言解說、ASCII 記憶體圖、gdb 實測輸出、
以及完整可執行的 \texttt{exploit.py} 逐行註解。
\end{minipage}
\end{center}

\vspace{0.6em}
\hrule height 0.4pt
\vspace{0.4em}

\renewcommand{\contentsname}{目錄}
\setcounter{tocdepth}{1}
\tableofcontents

\clearpage
}

\section{測試文件：CJK
排版驗證}\label{ux6e2cux8a66ux6587ux4ef6cjk-ux6392ux7248ux9a57ux8b49}

\subsection{一、段落與行內程式碼}\label{ux4e00ux6bb5ux843dux8207ux884cux5167ux7a0bux5f0fux78bc}

這是一段繁體中文測試，包含行內程式碼 \texttt{pop\ rdi;\ ret}、運算子 → ⇒
≥ ≤ ≠、星號 ★ ☆、 以及符號 ✔ ✗ ① ②
③。英文與中文混排應該要正確斷行，不應該出現方框或亂碼。

\begin{quote}
⚠️ 注意：這是一段引言區塊，用來測試引言樣式是否正確。
\end{quote}

\subsection{二、表格}\label{ux4e8cux8868ux683c}

{\def\LTcaptype{none} % do not increment counter
\begin{longtable}[]{@{}lll@{}}
\toprule\noalign{}
位址 & 指令 & 說明 \\
\midrule\noalign{}
\endhead
\bottomrule\noalign{}
\endlastfoot
\texttt{0x401069} & \texttt{syscall;\ ret} & 我們使用的
gadget（本題有兩個） \\
\texttt{0x40104c} & \texttt{syscall;\ ret} & leak 函式裡的，同樣可用 \\
\end{longtable}
}

\subsection{三、ASCII
圖表（含中文與製表符號）}\label{ux4e09ascii-ux5716ux8868ux542bux4e2dux6587ux8207ux88fdux8868ux7b26ux865f}

\begin{Shaded}
\begin{Highlighting}[]
\NormalTok{   高位址}
\NormalTok{   ┌──────────────────────────────┐}
\NormalTok{   │  return address（8 bytes）    │  ← 我們要覆寫的目標}
\NormalTok{   ├──────────────────────────────┤}
\NormalTok{   │  saved rbp（8 bytes）        │}
\NormalTok{   ├──────────────────────────────┤}
\NormalTok{   │  canary（8 bytes）           │  ← 最靠近 rbp，最先被撞到}
\NormalTok{   ├──────────────────────────────┤}
\NormalTok{   │  buf[0x30]                   │  ← rbp{-}0x40}
\NormalTok{   └──────────────────────────────┘}
\NormalTok{   低位址}
\end{Highlighting}
\end{Shaded}

\subsection{四、程式碼區塊}\label{ux56dbux7a0bux5f0fux78bcux5340ux584a}

\begin{Shaded}
\begin{Highlighting}[]
\ImportTok{from}\NormalTok{ pwn }\ImportTok{import} \OperatorTok{*}
\NormalTok{context.arch }\OperatorTok{=} \StringTok{\textquotesingle{}amd64\textquotesingle{}}
\NormalTok{p }\OperatorTok{=}\NormalTok{ process(}\StringTok{\textquotesingle{}./chal\textquotesingle{}}\NormalTok{)}
\NormalTok{payload  }\OperatorTok{=} \StringTok{b\textquotesingle{}A\textquotesingle{}} \OperatorTok{*} \BaseNTok{0x38} \OperatorTok{+}\NormalTok{ p64(canary) }\OperatorTok{+}\NormalTok{ p64(}\DecValTok{0}\NormalTok{) }\OperatorTok{+}\NormalTok{ p64(pop\_rdi) }\OperatorTok{+}\NormalTok{ p64(binsh) }\OperatorTok{+}\NormalTok{ p64(system)}
\NormalTok{p.sendline(payload)}
\BuiltInTok{print}\NormalTok{(}\StringTok{\textquotesingle{}這一行很長很長很長很長很長很長很長很長很長很長很長很長很長很長很長很長很長很長很長很長很長很長很長\textquotesingle{}}\NormalTok{)  }\CommentTok{\# 測試過長行的換行行為}
\end{Highlighting}
\end{Shaded}

\subsection{五、清單}\label{ux4e94ux6e05ux55ae}

\begin{enumerate}
\def\labelenumi{\arabic{enumi}.}
\tightlist
\item
  第一項
\item
  第二項

  \begin{itemize}
  \tightlist
  \item
    巢狀項目
  \item
    另一個巢狀項目
  \end{itemize}
\item
  第三項
\end{enumerate}

{\def\LTcaptype{none} % do not increment counter
\begin{longtable}[]{@{}ll@{}}
\toprule\noalign{}
項目 & 狀態 \\
\midrule\noalign{}
\endhead
\bottomrule\noalign{}
\endlastfoot
CJK 正文 & ✔ \\
等寬 CJK & ✔ \\
製表符號 & ✔ \\
\end{longtable}
}

\end{document}
